
% Chapter 2 File

\chapter{Technical Background}
\label{04_Technical_Background}
\thispagestyle{empty}

This chapter provides an overview of the features needed for the classification task, focusing on the selection and analysis of the most suitable dataset. Initial data analysis highlights a major difference between human-written and AI-generated content, emphasizing the importance of considering this when evaluating classifiers. The chapter then reviews the tools and technologies that support the development of the models, including programming languages, machine learning libraries, and computational resources.

\vspace{3mm}
Entire code base used in this research, including data processing scripts, model training code, and evaluation metrics, is publicly available on GitHub~\cite{neuman2025aicontentdetectors}. This repository contains all necessary resources for reproducing and extending this study.

\section{Data understanding}

To properly train effective classifiers, a well-balanced dataset with high-quality data is essential. Ideally, a training set is a representative subset of the population. In practice, we rarely achieve a perfectly representative set, as we often rely on random samples~\cite{borovicka2012selecting}. This challenge is even more prominent in the context of this paper, where the number of possible sentences is virtually infinite, let alone entire paragraphs or even longer pieces of~text. 

\vspace{3mm}

After carefully reviewing possible dataset options, OpenLLMText~\cite{chen2023openllmtext} stood out, due to its substantial size, diverse data sources, and the structured correlation between human-written and AI-generated texts. In this dataset, most prompts for generating AI text were human-written paragraphs, with the AI tasked to rephrase them. This set of data contains more than 350,000 text entries collected from 5 different sources:

\vspace{3mm}

\begin{itemize}
    \item  \textbf{Human-written texts}: around 70,000 entries randomly selected from the OpenWebText dataset, particularly from user-generated content from Reddit before 2019.
    \item \textbf{ChatGPT-generated texts}: around 70,000 entries created by GPT-3.5 Turbo, each being a paragraph-by-paragraph rephrasing of the human-written data.
    \item \textbf{LLaMA-7B-generated texts}: around 70,000 entries generated by the LLM Meta AI (LLaMA-7B) using similar paragraph-by-paragraph rephrasing techniques.
    \item \textbf{PaLM-generated texts}: around 70,000 entries produced by the Pathway Language Model (text-bison-001), also rephrasing human-written data.
    \item \textbf{GPT2-output dataset}: around 80,000 entries adapted from the GPT2-output dataset, released by OpenAI (GPT2-XL).
\end{itemize}

\vspace{3mm}

Google, the creator of PaLM, has announced that starting April 9, 2025, this model will not be accessible as it will be migrated to Gemini models. As a result, PaLM-generated data is excluded from our model training. Similarly, we omit using GPT-2 due to its lack of alignment with human-generated paragraphs and its outdated release (2019).

\vspace{3mm}

While GPT-3.5 Turbo and LLaMA-7B are not the most recent models -- the first one was released on March 14, 2023, and the second on July 18, 2023 -- they still offer high utility. Their alignment with human-written texts through paragraph-by-paragraph rephrasing is something unique that could not be found with newer models such as GPT-4 or Gemini.

\vspace{3mm}

Ensuring a balanced dataset for creating and evaluating our models involves using entries from the human-written subset while alternately selecting corresponding entries from GPT-3.5 Turbo and LLaMa. By doing so, we guarantee an equal class ratio of 1, resulting in a total of 134,886 entries. The dataset is then divided into three subsets~\cite{borovicka2012selecting}:

\vspace{3mm}

\begin{itemize}
    \item \textbf{Training set} (approximately 75\% of all entries) --  a set used for learning and estimating parameters of the model.
    \item \textbf{Validation set} (approximately 15\% of all entries) -- a set used to evaluate the model, usually for model selection.
    \item \textbf{Testing set} (approximately 10\% of all entries) -- a set of examples used to assess the predictive performance of the model.
\end{itemize}

\vspace{3mm}

Managing such a large dataset can be challenging during model training. To address this, we create sample datasets containing 50,000, 20,000, 5,000, 3,000 and 1,000 entries, each maintaining the 75-15-10 proportion for training, validation, and testing sets, respectively. These sample sets are used in the early stages of model development, allowing us to conserve time and computational resources. This approach enables us to focus on critical aspects such as feature engineering, hyperparameter selection, model architecture tuning, and evaluation metrics, which are essential for optimizing model performance.

\vspace{3mm}

% \noindent\textbf{\large Initial data exploration}
\subsection*{Initial data exploration}
\label{Data Exploration}

% \noindent During the first exploration of the data big difference in average text length between Human entries and those generated by AI can be observed: 

% \begin{figure}[h!]
%     \centering
%     \includegraphics[width=\textwidth]{Images/text_length_distribution.png}
%     \caption{Distribution of Text Length for Human and AI-Generated Text}
%     \label{fig:text_length_distribution}
% \end{figure}


% \noindent The Human distribution closely resembles Positively Skewed Unimodal Distribution with a peak around 1300 characters. Text lengths of AI-generated text are more unpredictable and resemble two binomial distributions superimposed on each other. Both distribution for ChatGPT and LLaMA are displayed in Figure \ref{fig:ai_text_length_distribution}.

% \begin{figure}[h!]
%     \centering
%     \includegraphics[width=\textwidth]{Images/ai_text_length_distribution.png}
%     \caption{Distribution of Text Length for AI-Generated Text}
%     \label{fig:ai_text_length_distribution}
% \end{figure}

% However, after carefully examining this we observe that ChatGPT distibution looks like a typical binomial distribution, expected result. However LLaMA texts are the reason for this strange distribution for all AI-texts. This distribution but itself causes strange 2 picks. The possible reason for such strange distribution is that for most of the entries LLaMA tries to generate rather short output -- this is a reason for a left-hand side pick, but for entries with a very long input that are fairly common in the dataset. 

%%%%%%%% XXX

During initial exploration of the data, we notice a significant difference in the average text length between human-written entries and those generated by AI. Texts created by humans had, on average, 572 words, while their rephrased by AI versions only had 304. During evaluation of the results, this difference must be carefully considered, because our classifiers should not rely on text length. Instead, they should focus on other features, such as language style, word choice, and syntactic structures, that differentiate human and AI-generated content. Histograms of these distributions are displayed in Figure~\ref{fig:text_length_distribution}.

\vspace{3mm}

\begin{figure}[h!]
    \centering
    \includegraphics[width=\textwidth]{Images/word_count.png}
    \caption{Distribution of text length for human and AI-generated text. The human-written text tends to be more consistent, while AI-generated text shows more variation.}
    \label{fig:text_length_distribution}
\end{figure}

The distribution for human text closely follows a positively skewed, unimodal pattern, peaking at around 230 words. On the other hand, the AI-generated text is more unpredictable, resembling two overlapping binomial distributions. These two distributions for ChatGPT and LLaMA are shown in Figure \ref{fig:ai_text_length_distribution}.

\vspace{3mm}

\begin{figure}[h!]
    \centering
    % \includegraphics[width=\textwidth]{Images/ai_word_count.png}
    \includegraphics[width=1\textwidth, height=2\textheight, keepaspectratio]{Images/ai_word_count.png}
    \caption{Distribution of text length for AI-generated text, showing two peaks for ChatGPT and LLaMA.}
    \label{fig:ai_text_length_distribution}
\end{figure}

Upon closer inspection, it is clear that the ChatGPT distribution follows a typical binomial pattern. However, the LLaMA-generated texts are responsible for the strange two-peak shape in the overall AI text distribution. This uncommon pattern is likely an outcome of LLaMA typically producing shorter outputs for most of the texts, creating the first peak on the left. But when the model is given long inputs, which are fairly common in the dataset, it tries to cut them during rephrasing due to model limitations.

%%%%%%% XXXX

% --------------------

\section{Tools and Technologies}

Python is selected as the sole programming language in this research due to its high utility in scripting and extensive access to libraries for data processing and machine learning. To address the challenges of managing long and complex scripts, Jupyter Notebooks are used. They offer features such as cell-based execution and a clear, interactive interface.

\vspace{3mm}

Regarding the libraries utilized, several key ones are worth highlighting:

\vspace{3mm}

\begin{itemize}
    \item \textbf{NumPy} -- a fundamental library for scientific computing with Python, providing support for arrays, matrices, and a wide range of mathematical functions for data manipulation and analysis.

    \item \textbf{Pandas} -- a library that allows easy manipulation of data through its primary structures, Series and DataFrames, which are equivalent to columns and tables, respectively. It enables efficient data loading, processing, and preparation for training or analysis of experimental results.

    \item \textbf{Scikit-learn} -- a widely used machine learning library in Python that provides a range of tools for model building, evaluation, and tuning. It is particularly useful during the creation of Logistic Regression and Support Vector Machines.

    \item \textbf{PyTorch} -- a high-performance deep learning library that offers an imperative and Pythonic programming style, making it easy to debug and maintain. PyTorch supports hardware accelerators such as GPUs, making it extremely well-suited for developing models discussed in this paper: Convolutional Neural Networks (CNNs), Bidirectional Long Short-Term Memory (BiLSTM), and Transformers.

    \item \textbf{Hugging Face} -- a platform and library focused on simplifying the implementation, training, and fine-tuning of transformer-based models such as BERT. The library offers pre-trained weights, tokenizers, and pipelines, which accelerate model deployment and fine-tuning. Hugging Face also provides a collaborative platform where users can share models, datasets, and other machine learning resources, creating a strong community-driven ecosystem.

    \item \textbf{TensorFlow} and \textbf{Keras} -- TensorFlow is an open-source framework for building and deploying machine learning models, with a focus on deep learning. It offers flexible and powerful tools for model training, evaluation, and deployment. Keras, now integrated with TensorFlow, provides a high-level API that additionally simplifies the process of designing, training, and experimenting with neural networks. Together, TensorFlow and Keras make it easier to work with complex models like BERT, especially when fine-tuning or deploying them in production environments.

    \item \textbf{Matplotlib} and \textbf{Seaborn} -- libraries used for data visualization. Matplotlib is a flexible tool for creating a variety of plots, while Seaborn builds on it, providing a higher-level interface for statistical graphics. Together, they enable clear and informative visualizations of model performance and data insights.

    \item \textbf{SHAP} and \textbf{LIME} -- both SHAP and LIME are widely used XAI libraries for model interpretability. SHAP provides consistent and mathematically grounded feature importance values, explaining the impact of each feature on predictions. LIME, on the other hand, approximates complex models locally with simpler, interpretable models, offering explanations for individual predictions. Both are valuable tools for making black-box models more transparent and understandable.
    
\end{itemize}

\vspace{3mm}

Since training models on a large dataset containing over 130,000 texts and aiming at high metric results requires significant computational power, Google Colaboratory and its paid Pro subscription are used. The subscription provides access to Tesla T4 and NVIDIA A100 GPUs -- two engines extensively used in this study. A total of 300 compute units were purchased during the research. Each Pro subscription costs \$12.29 per month and provides 100 compute units.

\vspace{3mm}

The Tesla T4 GPU offers 12.7 GB of system RAM and 15 GB of GPU VRAM. It consumes approximately 1.44 units per hour. With its relatively high computational power and low unit consumption, this engine is used during most of the training process and hyperparameter optimization.

\vspace{3mm}

The NVIDIA A100 GPU, on the other hand, provides 83.5 GB of system RAM and 40 GB of GPU VRAM. It consumes around 8.47 units per hour, enabling less than 12 hours of computation per subscription. While the A100 GPU is exceptionally powerful, with an estimated cost ranging from \$8,000 to \$10,000, it is reserved for final computations due to its higher unit consumption.

% For comparison, running one epoch of training on a locally available MacBook Pro 16-inch with an M1 chip (base configuration) took over [time]. In contrast, using Google Colab Pro with the Tesla T4 GPU reduced the time to [time]. It's important to note that even though Google Colab provides fast computational resources, it has limitations, such as a maximum 12-hour runtime per notebook session and memory limitations of 12,7 GB RAM and 15 GB of GPU VRAM. 

% \section{?  NOTES  ?}

% There are various methods for classifying texts. 

% During a learning process most learning algorithms use all instances from the given training set to estimate parameters of a model, but commonly lot of instances in the training set are useless. 

% especially in our example - numbers of sentences that can be created is infinite, not even mentioning whole paragraphs or longer texts.